% from ega2-1-fr.tex
\oldpage[II]{143}
\subsection{Critère valuatif de propreté}
\label{subsection:II.7.3-fr}

\begin{proposition}[7.3.1]
\label{II.7.3.1-fr}
Soient $A$ un anneau de valuation, $Y=\operatorname{Spec}(A)$, $b$ le point
générique de $Y$, $X$ un schéma intègre, $f:X\to Y$ un morphisme fermé tel
que $f^{-1}(b)$ se réduise à un point $x$ et que l'homomorphisme correspondant
$k(b)\to k(x)$ soit bijectif. Alors $f$ est un isomorphisme.
\end{proposition}

Comme $f$ est fermé et dominant, on a $f(X)=Y$~; il suffit
(\hyperref[I.4.2.2-fr]{I, 4.2.2}) de prouver que pour tout $y'\neq b$ dans $Y$,
il existe un \emph{seul} point $x'$ tel que $f(x')=y'$ et que l'homomorphisme
correspondant $\mathcal{O}_{y'}\to\mathcal{O}_{x'}$ soit bijectif, car $f$ sera
alors un homéomorphisme. Or, si $f(x')=y'$, $\mathcal{O}_{x'}$ est un anneau
local contenu dans $K=k(x)=k(b)$ et dominant $\mathcal{O}_{y'}$~; ce dernier
est l'anneau local $A_{y'}$, donc est un anneau de valuation
(\hyperref[II.7.1.1-fr]{7.1.1}) ayant $K$ pour corps des fractions. Par
ailleurs on a $\mathcal{O}_{x'}\neq K$ puisque $x'$ n'est pas le point
générique de $X$ (\hyperref[0.2.1.3-fr]{0, 2.1.3})~; on en conclut que
$\mathcal{O}_{x'}=\mathcal{O}_{y'}$. Comme $X$ est un schéma intègre, la
relation $\mathcal{O}_{x'}=\mathcal{O}_{x''}$ entraîne $x'=x''$
(\hyperref[I.8.2.2-fr]{I, 8.2.2}), ce qui achève la démonstration.

\begin{env}[7.3.2]
\label{II.7.3.2-fr}
Soient $A$ un anneau de valuation, $Y=\operatorname{Spec}(A)$, $b$ le point
générique de $Y$, de sorte que $\mathcal{O}_b=k(b)$ est égal à $K$, corps des
fractions de $A$~; soit $f:X\to Y$ un morphisme. On sait
(\hyperref[I.7.1.4-fr]{I, 7.1.4}) que les $Y$-sections \emph{rationnelles} de
$X$ sont en correspondance biunivoque avec les \emph{germes} de $Y$-sections
(définies dans des voisinages de $b$) au point $b$, d'où une application
canonique
\[
\tag{7.3.2.1}
\label{II.7.3.2.1-fr}
\Gamma_{\mathrm{rat}}(X/Y)\longrightarrow
\Gamma(f^{-1}(b)/\operatorname{Spec}(K))
\]
les éléments de $\Gamma(f^{-1}(b)/\operatorname{Spec}(K))$ s'identifiant
d'ailleurs, par définition (\hyperref[I.3.4.5-fr]{I, 3.4.5}) aux \emph{points
de $f^{-1}(b)=X\otimes_AK$ rationnels sur $K$}. Lorsque $f$ est \emph{séparé},
il résulte de (\hyperref[I.5.4.7-fr]{I, 5.4.7}) que l'application
(\hyperref[II.7.3.2.1-fr]{7.3.2.1}) est \emph{injective}, puisque $Y$ est un
schéma intègre.

Composant (\hyperref[II.7.3.2.1-fr]{7.3.2.1}) avec l'application canonique
$\Gamma(X/Y)\to\Gamma_{\mathrm{rat}}(X/Y)$
(\hyperref[I.7.1.2-fr]{I, 7.1.2}), on obtient une application canonique
\[
\tag{7.3.2.2}
\label{II.7.3.2.2-fr}
\Gamma(X/Y)\longrightarrow \Gamma(f^{-1}(b)/\operatorname{Spec}(K)).
\]
Lorsque $f$ est \emph{séparé}, cette application est encore \emph{injective}
(\hyperref[I.5.4.7-fr]{I, 5.4.7}).
\end{env}

\begin{proposition}[7.3.3]
\label{II.7.3.3-fr}
Soient $A$ un anneau de valuation de corps des fractions $K$,
$Y=\operatorname{Spec}(A)$, $b$ le point générique de $Y$, $f:X\to Y$ un
morphisme \emph{séparé} et \emph{fermé}. Alors l'application canonique
(\hyperref[II.7.3.2.2-fr]{7.3.2.2}) est bijective (ce qui revient à dire
qu'elle est \emph{surjective} et entraîne que les $Y$-sections rationnelles de
$X$ sont \emph{partout définies}).
\end{proposition}

Soit donc $x$ un point de $f^{-1}(b)$, \emph{rationnel sur $K$}. Comme $f$ est
séparé, il en est de même du morphisme
$f^{-1}(b)\to\operatorname{Spec}(K)$ correspondant à $f$
(\hyperref[I.5.5.1-fr]{I, 5.5.1, (iv)}), et toute section de $f^{-1}(b)$ étant
une immersion fermée (\hyperref[I.5.4.6-fr]{I, 5.4.6}), $\{x\}$ est fermé
\emph{dans $f^{-1}(b)$}. Considérons le sous-préschéma fermé réduit $X'$ de
$X$ ayant pour espace sous-jacent l'adhérence $\overline{\{x\}}$ de $\{x\}$
dans $X$. Il est clair que la restriction de $f$ à $X'$ vérifie les hypothèses
de (\hyperref[II.7.3.1-fr]{7.3.1}), donc est un \emph{isomorphisme} de $X'$ sur
$Y$, dont l'isomorphisme réciproque est la $Y$-section de $X$ cherchée.

\begin{env}[7.3.4]
\label{II.7.3.4-fr}
Pour énoncer les deux résultats suivants, nous nous servirons d'une
terminologie qui sera justifiée et discutée dans le chapitre IV~: si $F$ est
une partie
